\scanpage{12}
% archive sheet number: 11
% lines 1–2
\deleted{\uline{Corollaire}} \deleted{\uline{Proposition}} \inserted{\uline{Corollaire}} Supposons que l'on ait la condition \deleted{\unsure{(ii)}} suivante plus forte que (ii)

% lines 3–4
(ii$'$) \inserted{Les $\underline O_{X,x}$ noeth, et} Pour tt entier $n$, $\underline O_X[T_1,-T_n]$ est un faisceau d'anneaux cohérent.
\margin{?}

% lines 5
\inserted{(iii$'$) Tout $x\in X$ a un syst. fond. de voisinages \unsure{qu.-cpts|qui-compts} satisfont la condition (iii) pour le faisceau induit par $\underline O_X$ ....}

% lines 6–9
Alors la conclusion du lemme A est valable en y remplaçant $\underline O_X$ par $\underline O_X[T_1,-T_n]$, plus gén. par n'importe \unsure{quelle|quel} faisceau d'Algèbres $\underline B$ de présentation finie.

% lines 10–13
Soit $\underline B^p\mathrel{\overset{\inserted{u=(u_{ij})_{(i,j)\in P\times Q}}}{-}}B^q-M-0$ une suite exacte, \inserted{$\Sigma$} \deleted{\unsure{Alors on}} de $B$-Modules ($B=\underline O_X[T_1,-,T_n]$) Alors on a pour tt entier $m$ une suite exacte \inserted{$\Sigma_m$} de $\underline O_X$-Modules \inserted{cohérents}

% lines 14
\[
 \underline B_m^p - B_{N+m}^q\to M_m-0
\]

% lines 15
($N$ majorant des degrés des $u_{ij}$)

% lines 16
\deleted{\unsure{alors}} et $\Sigma$ est \uline{lim} des $\Sigma_m$. On aura

% lines 17
\[
 \Gamma B_m^p - \Gamma B_{\unsure{m+m|N+m|m+N}}^q - \Gamma M_m\to0
\]

% lines 18
d'où en passant à la limite

% lines 19
\[
 \Gamma B^p - \Gamma B^q - \Gamma M-0
\]

% lines 20
Ce qui implique $\alpha$. \deleted{\unsure{P}} Ceci prouve
